\documentclass{article}
\usepackage[english]{babel}
\usepackage{amsmath,amssymb,enumerate,theorem}

%%%%%%%%%% Start TeXmacs macros
\newcommand{\tmdummy}{$\mbox{}$}
\newcommand{\tmverbatim}[1]{\text{{\ttfamily{#1}}}}
\newenvironment{enumeratealpha}{\begin{enumerate}[a{\textup{)}}] }{\end{enumerate}}
\newenvironment{enumerateroman}{\begin{enumerate}[i.] }{\end{enumerate}}
\newenvironment{tmindent}{\begin{tmparmod}{1.5em}{0pt}{0pt}}{\end{tmparmod}}
\newenvironment{tmparmod}[3]{\begin{list}{}{\setlength{\topsep}{0pt}\setlength{\leftmargin}{#1}\setlength{\rightmargin}{#2}\setlength{\parindent}{#3}\setlength{\listparindent}{\parindent}\setlength{\itemindent}{\parindent}\setlength{\parsep}{\parskip}} \item[]}{\end{list}}
{\theorembodyfont{\rmfamily\small}\newtheorem{exercise}{Exercise}}
%%%%%%%%%% End TeXmacs macros

\begin{document}

\chapter*{1 Algebraic Background}

\begin{exercise}
  {\tmdummy}
  
  \begin{tmindent}
    Domain $\mathbb{Z}$ is not a field.
  \end{tmindent}
\end{exercise}

\begin{exercise}
  {\tmdummy}
  
  \begin{enumeratealpha}
    \item irreducible
    
    \item $(x + 13) (x - 13)$
    
    \item $(x + 1) (x^2 + 1)$
    
    \item Irreducibility can be obtained by computing the cases of $x \in \{ -
    3, - 2, - 1, 0, 1 \}$ and checking \ monotonicity.
  \end{enumeratealpha}
\end{exercise}

\begin{exercise}
  {\tmdummy}
  
  \begin{tmindent}
    $x^2 - 1 = (x + 1) (x - 1)$
  \end{tmindent}
\end{exercise}

\begin{exercise}
  {\tmdummy}
  
  \begin{tmindent}
    Yes.
    
    In fact, for any field $F$ and $f (x) \in F [x]$. We'll prove $\exists K /
    F, s.t.f$ have multiple roots in $K \Longleftrightarrow f, f'$ are not
    relatively prime: Let $\alpha$ be a root of $f$, $f (x) = (x - \alpha) g
    (x)$. Now $f' (x) = g (x) + (x - \alpha) g' (x)$. Substituing $x =
    \alpha$, we'll see that $f' (\alpha) = 0 \Leftrightarrow g (\alpha) = 0
    \Leftrightarrow f (x)$ has multiple roots.
  \end{tmindent}
\end{exercise}

\begin{exercise}
  {\tmdummy}
  
  \begin{enumerateroman}
    \item $x^4 + 1$
    
    \item $x^4 - 2 x^2 + 9$
    
    \item $x^3 - 6 x^2 + 12 x - 9$
  \end{enumerateroman}
\end{exercise}

\begin{exercise}
  {\tmdummy}
  
  \begin{enumeratealpha}
    \item 2
    
    \item 1
    
    \item 4
    
    \item 3
    
    \item infinite
  \end{enumeratealpha}
\end{exercise}

\begin{exercise}
  {\tmdummy}
  
  \begin{tmindent}
    It is an algebraic extension for sure since its generators are all
    algebraic elements.
    
    However, it's not finite since the minimal polynomial of each of the
    generators can be arbitrarily high.
  \end{tmindent}
\end{exercise}

\begin{exercise}
  {\tmdummy}
  
  \begin{enumeratealpha}
    \item $s_1^2 - 2 s_2$
    
    \item $s_1^3 - 3 s_1 s_2$
    
    \item It's not symmetric.
    
    \item It's not symmetric.
  \end{enumeratealpha}
\end{exercise}

\begin{exercise}
  {\tmdummy}
  
  \begin{enumerateroman}
    \item We can prove $\Delta$ is antisymmetic by comparing the definition of
    even permutations and components of $\Delta$.
    
    \item First, we'll prove that for $f (x_1, x_2, \ldots, x_n) \in
    \mathbb{Z} [x_1, x_2, \ldots, x_n]$ which is antisymmetric, $\Delta \mid f
    (x_1, x_2, \ldots, x_n)$: To prove this, we only have to prove that
    $\exists i, j \in [1, n] \cap \mathbb{Z}$, let's say $i = 1, j = 2$, $(x_i
    - x_j) \mid f (x_1, x_2, \ldots, x_n)$. Notice $f (x_1, x_2, \ldots, x_n)
    \in (\mathbb{Z} [x_3, x_4, \ldots, x_n]) [x_1, x_2]$, we can have
    \[ f (x_1, x_2, \ldots, x_n) = \Large{\Sigma}_{i, j \in [0, n] \cap
       \mathbb{Z}} a_{i j} x_1^i x_2^j \]
    where $a_{i j} \in \mathbb{Z} [x_3, x_4, \ldots, x_n]$. And because $f
    (x_1, x_2, \ldots, x_n) = - f (x_2, x_1, \ldots, x_n)$, $a_{i j} = - a_{j
    i}$, indicating $(x_1 - x_2) \mid f (x_1, x_2, \ldots, x_n)$. Thus,
    $\Delta \mid f (x_1, x_2, \ldots, x_n)$.
    
    Now, following form the hint, we consider $f / \Delta$. It is a symmetic
    polynomial, so it was consist of elementary symmetric polynomials.
  \end{enumerateroman}
\end{exercise}

\begin{exercise}
  \tmverbatim{I'm not quite sure about b) and d)}
  \begin{enumeratealpha}
    \item 42
    
    \item $\left[\begin{array}{ccc}
      1 & 3 & 5\\
      2 & - 4 & \\
      7 & 2 & 9
    \end{array}\right] \rightarrow \left[\begin{array}{ccc}
      1 & 3 & 5\\
      & 10 & - 10\\
      & 11 & 14
    \end{array}\right] \rightarrow \left[\begin{array}{ccc}
      1 & 3 & 5\\
      & 1 & 24\\
      & 10 & - 10
    \end{array}\right] \rightarrow \left[\begin{array}{ccc}
      1 & 3 & 5\\
      & 1 & 24\\
      &  & - 250
    \end{array}\right]$, the answer is 250.
    
    \item infinite
    
    \item $\left[\begin{array}{ccc}
      41 & 32 & - 999\\
      & 16 & 3\\
      & 2 & 111
    \end{array}\right] \rightarrow \left[\begin{array}{ccc}
      41 & 50 & 0\\
      & 16 & 3\\
      & 2 & 111
    \end{array}\right] \rightarrow \left[\begin{array}{ccc}
      41 & 2 & - 9\\
      & 16 & 3\\
      & 2 & 111
    \end{array}\right] \rightarrow \left[\begin{array}{ccc}
      41 & 2 & - 9\\
      & 2 & 111\\
      &  & - 885
    \end{array}\right]$, the answer is $412885 = 72, 570$.
    
    \item infinite
  \end{enumeratealpha}
\end{exercise}

\begin{exercise}
  {\tmdummy}
  
  \begin{enumerateroman}
    \item The proff can be obtained by checking the definitions.
    
    \item The proff can be obtained by checking the definitions.
    
    \item I suppose it can't because we no longer have $1 m = m$, which means
    that the values of mutiplication $km$ are no longer distributed to the
    vectors in $M$.
  \end{enumerateroman}
\end{exercise}

\begin{exercise}
  {\tmdummy}
  
  \begin{tmindent}
    They are $n \mathbb{Z}$.
  \end{tmindent}
\end{exercise}

\begin{exercise}
  {\tmdummy}
  
  \begin{tmindent}
    No. For example, $\mathbb{Z}$ has infinite numbers of submodules $n
    \mathbb{Z}$.
    
    Such condition can be obtained by have $R$ as a finite ring.
  \end{tmindent}
\end{exercise}

\begin{exercise}
  \ 
\end{exercise}

\begin{exercise}
  \ 
\end{exercise}

\end{document}
